\documentclass{article}
\usepackage{amsmath,amssymb}
\usepackage[margin=1in]{geometry}
\usepackage{hyperref}

\title{CSDR--Based EFT Positivity Bounds: Simplification Record}
\author{sdpb pipeline}
\date{\today}

\begin{document}
\maketitle

\section*{Conventions and Basis}

\paragraph{Kinematic variables.}
Let $s,t,u$ be the standard Mandelstam variables with $s+t+u = 4m^2 \equiv \mu$.
Define
\[
  S_1 = s - \tfrac{\mu}{3},\quad S_2 = t - \tfrac{\mu}{3},\quad S_3 = u-\tfrac{\mu}{3},
\]
so that $S_1+S_2+S_3=0$.  The amplitude is expanded in the crossing-symmetric basis
\[
  \mathcal{M} = \sum_{p,q\ge0} W_{p,q}\, x^p\, y^q,
  \qquad x = -(S_1 S_2+S_2 S_3+S_3 S_1),\quad y = -S_1 S_2 S_3.
\]
The $W_{p,q}$ are the Wilson coefficients of the EFT amplitude.

\textbf{Important:} $S_1,S_2,S_3$ are \emph{not} Mandelstam variables; they are
Mandelstam variables minus $\mu/3$.  In particular $S_1+S_2+S_3=0$.

\paragraph{Classification of $W_{p,q}$.}
Following the CSDR paper and the user's notes:
\begin{itemize}
  \item $p\ge0$ (first index $\ge0$): \textbf{objectives} (EFT coefficients; $W_{p,q}\ge0$
        from unitarity in the CSDR scheme).
  \item $p<0$ (first index $<0$): \textbf{null constraints} ($W_{p,q}=0$
        by the CSDR crossing symmetry).
\end{itemize}
In CSDR notation this is $W_{n-m,m}$ with $n-m\ge0$ (objectives) or $n-m<0$ (null).

\paragraph{Spacetime dimension.}
$d$ is a free parameter throughout.  Define
\[
  \alpha = \frac{d-3}{2},
\]
so that the Gegenbauer parameter satisfies $2\alpha = d-3$ and the spin measure
$(2\ell+2\alpha) = (2\ell+d-3)$ matches the user's eq.~(1).

\section{Heavy-Average Formulation (User Notes eq.~(2))}

\paragraph{Starting point.}
From the CSDR dispersion relation (Sinha--Zahed eq.~(4)):
\[
  W_{n-m,\,m}
  = \int_{\delta_0}^\infty \frac{ds_1}{s_1}\,\Phi(s_1)
    \sum_{\ell=0,\text{even}}^\infty (2\ell+2\alpha)\,a_\ell(s_1)\,B_{\ell}^{(n,m)}(s_1).
\]

\paragraph{Step 1: large-$s_1$ limit of $\Phi$.}
$\Phi(s_1) = s_1^{(4-d)/2}\,\Psi(\alpha)$ at large $s_1$, where $\Psi(\alpha)>0$ is a
positive constant.  Since $\Psi(\alpha)$ is an overall scale, we set $\Psi(\alpha)=1$.

\paragraph{Step 2: large-$s_1$ limit of the kernel.}
The CSDR kernel $B_\ell^{(n,m)}(s_1)$ has the large-$s_1$ expansion
\[
  B_\ell^{(n,m)}(s_1)
  \;\xrightarrow{s_1\to\infty}\;
  \frac{1}{\pi}\,C_\ell^{(\alpha)}(1)\,\frac{D_{\ell,\alpha}^{(n,m)}}{s_1^{2n+m}},
\]
where the coefficient $D_{\ell,\alpha}^{(n,m)}$ is given by CSDR eq.~(11) (valid for
$m>n\ge1$):
\[
  D_{\ell,\alpha}^{(n,m)}
  = \Gamma(m-n)\sum_{j=n}^{m}
    \frac{(-4)^j\,\bigl(-\tfrac{\ell}{2}\bigr)_j\,\bigl(\alpha+\tfrac{\ell}{2}\bigr)_j
          \,(3j-m-2n)}
         {j!\,(m-j)!\,(j-n)!\,\bigl(\alpha+\tfrac12\bigr)_j}.
\]

\paragraph{Step 3: Identify with the heavy average.}
The Extremal EFT heavy average (eq.~(1) of user notes) is
\[
  \langle F\rangle
  = \sum_\ell n_\ell^{(d)}
    \int_{\delta_0}^\infty \frac{ds_1}{s_1}\,\frac{s_1^{4-d}}{\pi}\,\rho_\ell(s_1)\,F(s_1,\ell),
\]
with $n_\ell^{(d)} = \dfrac{(4\pi)^{d/2}(d+2\ell-3)\,\Gamma(d+\ell-3)}{\pi\,\Gamma\!\bigl(\tfrac{d-2}{2}\bigr)\,\Gamma(\ell+1)}$.

The CSDR partial-wave coefficient satisfies $a_\ell = n_\ell^{(d)}\,\rho_\ell/\pi$ (up to
normalisation).  Substituting and using $(2\ell+2\alpha)=(2\ell+d-3)$:

\[
  W_{n-m,\,m}
  = \Bigl\langle
      \frac{D_{\ell,\alpha}^{(n,m)}\,C_\ell^{(\alpha)}(1)\,(2\ell+d-3)}{n_\ell^{(d)}\,s_1^{2n+m}}
    \Bigr\rangle.
\]

\paragraph{Step 4: Cancel $n_\ell^{(d)}$.}
The factor $n_\ell^{(d)}$ appears in both the measure $\langle\cdot\rangle$ and the
integrand.  It cancels, so the SDPB polynomial blocks do \emph{not} need it.  The
result is user notes \textbf{eq.~(2)}:
\begin{equation}\label{eq:W_kernel}
  \boxed{
    W_{n-m,\,m}
    = \Bigl\langle
        \frac{D_{\ell,\alpha}^{(n,m)}\,C_\ell^{(\alpha)}(1)\,(2\ell+d-3)}{s_1^{2n+m}}
      \Bigr\rangle
  }
\end{equation}
valid for $m>n\ge1$ (null constraints).

\paragraph{Objectives ($n\ge m\ge0$, $n\ge1$).}
For $n\ge m$, $D^{(n,m)}=0$ (empty sum in eq.~(11)).  The forward-scattering spectral
function gives a positive-definite kernel:
\[
  W_{n-m,\,m}\;\text{(objective)}
  = \Bigl\langle
      \frac{C_\ell^{(\alpha)}(1)\,(2\ell+d-3)}{s_1^{2n+m}}
    \Bigr\rangle \ge0.
\]
This is consistent with unitarity ($a_\ell\ge0$) in the CSDR scheme.

\section{SDPB Polynomial Matrix Program}

\paragraph{Decision variables.}
\begin{itemize}
  \item $z_{p,q}$ for each objective pair $(p,q)$ with $p\ge1$, $q\ge0$, $2p+3q\le K$.
  \item $c_{n,m}$ for each null pair $(n,m)$ with $m>n\ge1$, $2n+m\le K$.
\end{itemize}

\paragraph{Spectral polynomial per spin $\ell$.}
Setting $s_1 = M^2(1+x)$ (threshold at $x=0$) and clearing the common denominator
$(1+x)^K$, the SDPB polynomial for even spin $\ell$ is:
\begin{align}
  P^{(\ell)}(x)
  &= \sum_{(p,q)\in\text{obj}}
       z_{p,q}\underbrace{C_\ell^{(\alpha)}(1)(2\ell+d-3)}_{\kappa^{\rm obj}_{p,q,\ell}}
       (1+x)^{K-2p-3q} \nonumber\\
  &\quad+ \sum_{(n,m)\in\text{null}}
       c_{n,m}\underbrace{D_{\ell,\alpha}^{(n,m)}C_\ell^{(\alpha)}(1)(2\ell+d-3)}_{\kappa^{\rm null}_{n,m,\ell}}
       (1+x)^{K-2n-m}.
\end{align}

\paragraph{SDPB problem.}
\begin{alignat}{2}
  &\text{maximize}  &&\quad z_{p_0,q_0} \nonumber\\
  &\text{subject to}&&\quad z_{p_1,q_1} = 1 \quad\text{(normalization)}\nonumber\\
  &                 &&\quad P^{(\ell)}(x)\ge0\quad\forall\,\ell\in\{0,2,\ldots,\ell_{\max}\},\,\forall\, x\ge0.
\end{alignat}

\paragraph{DampedRational prefactor.}
Each spin block has prefactor $e^{-x}/(1+x)^K$ (pole at $x=-1$, multiplicity $K$).

\section{Simplification Record for the Code}

The following steps are performed by the code in order.

\begin{enumerate}
  \item \texttt{csdr.alpha\_from\_d(d)} computes $\alpha=(d-3)/2$ as an exact
        \texttt{Fraction}.

  \item \texttt{csdr.gegenbauer\_at\_one(ell,alpha)} computes
        $C_\ell^{(\alpha)}(1) = (2\alpha)_\ell / \ell!$ using Pochhammer symbols.

  \item \texttt{csdr.D\_coeff(n,m,ell,alpha)} computes $D_{\ell,\alpha}^{(n,m)}$
        by evaluating the sum in CSDR eq.~(11) term by term as exact fractions.
        Verified against the closed form
        $D^{(1,2)} = 2\ell(\ell+2\alpha)(-11-10\alpha+2\ell(\ell+2\alpha))/((2\alpha+1)(2\alpha+3))$
        for $\ell=0,2,4,6,8$.

  \item \texttt{csdr.null\_kernel\_coeff(n,m,ell,d)} $= D \cdot C_\ell^{(\alpha)}(1) \cdot (2\ell+d-3)$
        (product of three exact fractions).

  \item \texttt{csdr.obj\_kernel\_coeff(p,q,ell,d)} $= C_\ell^{(\alpha)}(1) \cdot (2\ell+d-3)$.

  \item \texttt{csdr.poly\_expand\_1px(k)} expands $(1+x)^k$ using binomial coefficients.

  \item \texttt{pmp\_generator.\_csdr\_spin\_block} multiplies each kernel by the
        appropriate $(1+x)^{K-\text{power}}$ polynomial and pads to length $K+1$.

  \item \texttt{pmp\_generator.generate\_csdr\_pmp} assembles the objective vector,
        normalization vector, and all spin blocks into the SDPB PMP JSON.
\end{enumerate}

\section{Self-Consistency Checks}

The function \texttt{run\_bounds.run\_csdr\_checks(d)} verifies:
\begin{itemize}
  \item $D^{(1,2)}_{\ell,\alpha}$ matches the closed form for $\ell=0,2,4,6,8$.
  \item $C_\ell^{(\alpha)}(1) = 1$ for all $\ell$ when $d=4$ ($\alpha=1/2$).
  \item The null kernel $\kappa^{\rm null}_{1,2,2}$ is non-zero.
  \item The objective kernel $\kappa^{\rm obj}_{1,0,\ell}$ is positive for $\ell=0,2$.
  \item Pair enumeration is non-empty for $K=8$.
\end{itemize}

\section{Key Differences from the Extremal EFT Paper}

The Extremal EFT paper (Caron-Huot \& Duong) uses a fixed-$t$ dispersion relation
with a different subtraction scheme.  Its spectral functions are
\[
  v^{\rm EFT}_{p,q,\ell}(x) \propto
  \bigl[\text{Taylor coeff.\ of } C_\ell\bigr]_q \cdot x^{q_{\max}-q}(x+4m^2)^{p_{\max}-p},
\]
which is \emph{not} the same as the CSDR kernel $C_\ell^{(\alpha)}(1)(2\ell+d-3)/s_1^{2p+3q}$.

The two formulations use different subtraction schemes and will in general produce
different numerical bounds.  Getting the same bounds from both is a sign of error, not
correctness.

\end{document}
